% Generated draft: requires mathematical and source review.
\documentclass[11pt]{article}
\usepackage[margin=25mm]{geometry}
\usepackage{fontspec}
\setmainfont{Latin Modern Roman}
\usepackage{amsmath,amssymb,mathtools,unicode-math}
\setmathfont{Latin Modern Math}
\usepackage{enumitem,needspace,xurl,hyperref}
\setlength{\parindent}{0pt}
\setlength{\parskip}{5pt}
\setlength{\emergencystretch}{3em}
\newcommand{\N}{\mathbb{N}}
\newcommand{\Z}{\mathbb{Z}}
\newcommand{\Q}{\mathbb{Q}}
\newcommand{\R}{\mathbb{R}}
\newcommand{\C}{\mathbb{C}}
\newcommand{\F}{\mathbb{F}}
\newcommand{\A}{\mathbb{A}}
\newcommand{\PP}{\mathbb{P}}
\newcommand{\E}{\mathbb{E}}
\newcommand{\eps}{\varepsilon}
\newcommand{\dd}{\,\mathrm{d}}
\begin{document}
\section*{UnsolvedMath ID: 30004934}
\textbf{Source problem code:} catalogue entry
\subsection{Definitions and mathematical statement}
Let \(k\geq 2\) be an integer, let \(G\) be either \(\R^k\) or \(\Z^k\), and let \(M\) be a compact smooth manifold. A \(C^\infty\) action is a smooth map \(\alpha:G\times M\to M\) satisfying \(\alpha(0,x)=x\) and \(\alpha(g+h,x)=\alpha(g,\alpha(h,x))\). Write \(\alpha_g(x)=\alpha(g,x)\).

The supplied report calls such an action \emph{Anosov} if there exists \(g\in G\) such that \(\alpha_g\) is normally hyperbolic with respect to the orbit foliation \(\mathcal O\) of \(G\). Normal hyperbolicity here refers to invariant contracting and expanding directions transverse to the orbit directions, with their contraction and expansion dominating the dynamics tangent to the orbits. The supplied passage stops at the beginning of the tangent-splitting definition, so its complete technical conventions are unavailable.

For two actions \(\alpha\) and \(\beta\) of the same group, a \(C^\infty\) conjugacy is a smooth diffeomorphism \(h\) satisfying
\[
 h\circ\alpha_g=\beta_g\circ h\qquad\text{for every }g\in G.
\]
A factor relation is expressed by an equivariant surjection \(\pi:M\to N\), satisfying \(\pi\circ\alpha_g=\gamma_g\circ\pi\) for a factor action \(\gamma\). This general description does not specify the additional regularity or nontriviality requirements intended for a \(C^\infty\) rank-one factor.

\textbf{Katok--Spatzier conjecture, as supplied.} Every \(C^\infty\) Anosov action of \(\R^k\) or \(\Z^k\), with \(k\geq2\), on a compact manifold, having no \(C^\infty\) rank-one factors, is \(C^\infty\)-conjugate to an algebraic action after passage to finite covers.

\textbf{Unresolved qualifications.} The supplied record does not define precisely which factors count as rank-one factors, which algebraic actions are allowed, or what operations are included in passage to finite covers. In particular, it does not specify whether a finite-index restriction of a \(\Z^k\) action is permitted when forming a lift. It also does not state local freeness or boundary conventions for \(M\), despite using an orbit foliation. These conventions must be recovered from the cited source before this can be treated as a fully specified classification assertion; no additional hypotheses or choices are imposed here.
\subsection{Short English statement}
The supplied Katok--Spatzier conjecture asks whether every smooth higher-rank Anosov action on a compact manifold, without smooth rank-one factors, becomes smoothly conjugate to an algebraic action after passing to finite covers. The precise meanings of rank-one factor, algebraic action, and the permitted covering operations remain unspecified in the supplied excerpt.
\subsection{Sources}
\begin{itemize}
\item Supplied local UnsolvedMath registry record: permanent ID \(30004934\), title \emph{Algebraic Rigidity of Higher-Rank Anosov Actions}, Oberwolfach Reports Open Problems collection, exact source problem code \texttt{catalogue entry}. The record supplies the labels \texttt{open} and \texttt{corrected\_verified}; these are registry assertions, not independent verification in this draft.
\item \emph{Mini-Workshop: \((Anosov)^3\) (2022)}, Oberwolfach Reports, pp.\ 3067--3099, DOI \texttt{10.4171/owr/2021/56}: \url{https://doi.org/10.4171/owr/2021/56}. The supplied excerpt is headed \emph{Mini Course on Higher Rank Anosov Actions}, by Ralf Spatzier, and attributes the conjecture to Katok--Spatzier. It says the problem was posed as a question in reference [3] and as a conjecture in reference [6, Conjecture 16.8]; the corresponding bibliography entries are not supplied. The registry gives proposed year 2021, while its report citation includes 2022.
\item Supplied literature reference: Spatzier--Vinhage, \emph{Cartan actions and classification}, \url{https://doi.org/10.1090/jams/1033}.
\item Supplied literature reference, labeled ``current survey of the conjecture'': \url{https://sites.lsa.umich.edu/spatzier/wp-content/uploads/sites/1081/2023/07/Katoks-work-Chapter4-final.pdf}.
\item Literature assessment dated exactly \texttt{2026-08-22}: the supplied curation labels the general Katok--Spatzier classification open and reports major rigidity results for important subclasses, including totally Cartan actions, without claiming the universal conjecture. This is supplied dated curation derived from the OWR record, not an independently verified literature assessment. The record also reports that its clean statement was checked against the official report; that source-verification assertion has not been independently checked here.
\end{itemize}
\end{document}
